% Bagian: computability
% Bab: computability-theory
% Subbagian: equiv-ce-defs

\documentclass[../../../include/open-logic-section]{subfiles}

\begin{document}

\olfileid[id]{cmp}{thy}{eqc}

\olsection[Definisi Himpunan C.E.]{Definisi Ekuivalen Himpunan
  Terenumerasi Secara Komputabel}


Pernyataan berikut memberikan sejumlah pencirian ekuivalen yang penting bagi
makna terenumerasi secara komputabel.

\begin{thm}
\ollabel{thm:ce-equiv}
Misalkan $S$ suatu himpunan bilangan asli. Maka pernyataan-pernyataan berikut
ekuivalen:
\begin{enumerate}
\item\ollabel{case:ce} $S$ terenumerasi secara komputabel.
\item\ollabel{case:ran-pc} $S$ merupakan daerah hasil suatu fungsi dapat
  dihitung \emph{parsial}.
\item\ollabel{case:ran-prim} $S$ kosong atau merupakan daerah hasil suatu
  fungsi rekursif primitif.
\item\ollabel{case:ce-domain} $S$ merupakan \emph{domain} suatu fungsi dapat
  dihitung parsial.
\end{enumerate}
\end{thm}

\begin{explain}
Tiga klausa pertama menyatakan bahwa kita dapat secara ekuivalen menganggap
setiap himpunan tak kosong yang terenumerasi secara komputabel dienumerasi oleh
suatu fungsi dapat dihitung, fungsi dapat dihitung parsial, atau fungsi
rekursif primitif. Klausa keempat memberi tahu kita bahwa jika $S$
terenumerasi secara komputabel, maka untuk suatu indeks~$e$,
\[
S = \Setabs{x}{\cfind{e}(x) \fdefined}.
\]
Dengan kata lain, $S$ adalah himpunan masukan yang membuat komputasi
$\cfind{e}$ berhenti. Karena alasan itu, himpunan terenumerasi secara komputabel
kadang disebut \emph{semidapat diputuskan}: jika suatu bilangan berada dalam
himpunan tersebut, pada akhirnya Anda memperoleh jawaban ``ya,'' tetapi jika
tidak berada di dalamnya, Anda tidak pernah memperoleh jawaban ``tidak''!{}
\end{explain}

\begin{proof}
Karena setiap fungsi rekursif primitif dapat dihitung dan setiap fungsi dapat
dihitung juga dapat dihitung secara parsial, \olref{case:ran-prim}
mengimplikasikan \olref{case:ce}, sedangkan \olref{case:ce}
mengimplikasikan~\olref{case:ran-pc}. (Perhatikan bahwa jika $S$ kosong,
$S$ merupakan daerah hasil fungsi dapat dihitung parsial yang tidak
terdefinisi di mana pun.) Jika kita menunjukkan bahwa \olref{case:ran-pc}
mengimplikasikan \olref{case:ran-prim}, kita telah menunjukkan bahwa tiga
klausa pertama ekuivalen.

Jadi, andaikan $S$ merupakan daerah hasil fungsi dapat dihitung parsial
$\cfind{e}$. Jika $S$ kosong, kita selesai. Jika tidak, misalkan $a$ sembarang
unsur~$S$. Berdasarkan teorema bentuk normal Kleene, kita dapat menulis
\[
\cfind{e}(x) \simeq U(\umin{s}{T(e, x, s)}).
\]
Khususnya, $\cfind{e}(x) \fdefined$ dan $= y$ jika dan hanya jika terdapat
suatu $s$ sedemikian sehingga $T(e, x, s)$ dan $U(s) = y$. Definisikan $f(z)$
dengan
\[
f(z) = \begin{cases}
  U((z)_1) & \text{jika $T(e, (z)_0, (z)_1)$} \\
  a        & \text{selain itu.}
\end{cases}
\]
Maka $f$ bersifat rekursif primitif karena $T$ dan $U$ bersifat demikian.
\iftag{TMs}{Dalam bahasa mesin Turing, jika $z$ mengodekan pasangan
  $\tuple{(z)_0, (z)_1}$ sedemikian sehingga $(z)_1$ merupakan komputasi
  berhenti mesin~$M_e$ pada masukan $(z)_0$, maka $f$ mengembalikan keluaran
  komputasi itu; jika tidak, $f$ mengembalikan~$a$.}\par

Kita harus menunjukkan bahwa $S$ merupakan daerah hasil~$f$, yakni bagi setiap
bilangan asli~$y$, $y \in S$ jika dan hanya jika $y$ berada dalam daerah
hasil~$f$. Untuk arah maju, andaikan $y \in S$. Maka $y$ berada dalam daerah
hasil $\cfind{e}$, sehingga untuk suatu $x$ dan~$s$, berlaku $T(e,x,s)$ dan
$U(s) = y$. Namun, dalam hal itu $y = f(\tuple{x,s})$. Sebaliknya, andaikan
$y$ berada dalam daerah hasil~$f$. Maka $y = a$, atau untuk suatu~$z$ berlaku
$T(e,(z)_0,(z)_1)$ dan $U((z)_1) = y$. Karena dalam kasus terakhir
$\cfind{e}((z)_0) \fdefined = y$, dalam kedua kasus $y$ berada dalam~$S$.

(Notasi $\cfind{e}(x) \fdefined = y$ berarti ``$\cfind{e}(x)$ terdefinisi dan
sama dengan $y$.'' Kita juga dapat menggunakan $\cfind{e}(x) = y$, tetapi
panah tambahan itu kadang membantu mengingatkan bahwa kita sedang berurusan
dengan suatu fungsi parsial.)

Untuk menyelesaikan bukti \olref{thm:ce-equiv}, cukup ditunjukkan bahwa
\olref{case:ce} dan~\olref{case:ce-domain} ekuivalen. Pertama, mari kita
tunjukkan bahwa \olref{case:ce} mengimplikasikan~\olref{case:ce-domain}.
Andaikan $S$ merupakan daerah hasil suatu fungsi dapat dihitung~$f$, yakni
\[
S = \Setabs{y}{\text{untuk suatu $x$, } f(x) = y}.
\]
Misalkan
\[
g(y) = \umin{x}{(f(x) = y)}.
\]
Maka $g$ merupakan fungsi dapat dihitung parsial, dan $g(y)$ terdefinisi jika
dan hanya jika, untuk suatu~$x$, $f(x) = y$. Dengan kata lain, domain~$g$
merupakan daerah hasil~$f$. \iftag{TMs}{Dalam bahasa mesin Turing: jika
diberikan mesin Turing~$F$ yang mengenumerasi unsur-unsur~$S$, misalkan $G$
adalah mesin Turing yang memutuskan-sebagian $S$ dengan menelusuri keluaran~$F$
untuk melihat apakah suatu unsur yang diberikan berada dalam himpunan itu; $G$
berhenti jika unsur tersebut ditemukan dan terus mencari selamanya jika tidak
ditemukan.}{}

Terakhir, untuk menunjukkan bahwa \olref{case:ce-domain}
mengimplikasikan~\olref{case:ce}, andaikan $S$ merupakan domain fungsi dapat
dihitung parsial~$\cfind{e}$, yakni
\[
S = \Setabs{x}{\cfind{e}(x) \fdefined}.
\]
Jika $S$ kosong, kita selesai; jika tidak, misalkan $a$ sembarang unsur~$S$.
Definisikan $f$ dengan
\[
f(z) = \begin{cases}
(z)_0 & \text{jika $T(e,(z)_0,(z)_1)$} \\
a & \text{selain itu.}
\end{cases}
\]
Maka, seperti di atas, suatu bilangan $x$ berada dalam daerah hasil~$f$ jika
dan hanya jika $\cfind{e}(x) \fdefined$, yakni jika dan hanya jika $x \in S$.
\iftag{TMs}{Dalam bahasa mesin Turing: jika diberikan mesin $M_e$ yang
memutuskan-sebagian $S$, enumerasilah unsur-unsur $S$ dengan menjalankan semua
komputasi mesin Turing yang mungkin dan mengembalikan masukan-masukan yang
bersesuaian dengan komputasi yang berhenti.}{}
\end{proof}

Klausa~\olref{case:ce-domain} dalam \olref{thm:ce-equiv} memberi kita cara yang
mudah untuk mengenumerasi himpunan terenumerasi secara komputabel: untuk
setiap~$e$, misalkan $W_e$ menyatakan domain $\cfind{e}$, yakni
\[
W_e = \Setabs{x}{\cfind{e}(x) \fdefined}.
\]
Maka jika $A$ merupakan sembarang himpunan terenumerasi secara komputabel,
$A = W_e$ untuk suatu~$e$.

Teorema berikut memberikan satu lagi pencirian himpunan terenumerasi secara
komputabel.

\begin{thm}
\ollabel{thm:exists-char}
Suatu himpunan $S$ terenumerasi secara komputabel jika dan hanya jika terdapat
relasi dapat dihitung $R(x,y)$ sedemikian sehingga
\[
S = \Setabs{ x }{ \lexists[y][R(x,y)] }.
\]
\end{thm}

\begin{proof}
Untuk arah maju, andaikan $S$ terenumerasi secara komputabel. Maka untuk suatu
$e$, $S = W_e$. Untuk nilai~$e$ ini kita dapat menulis $S$ sebagai
\[
S = \Setabs{ x }{ \lexists[y][T(e, x, y)] }.
\]
Untuk arah sebaliknya, andaikan $S = \Setabs{ x }{
  \lexists[y][R(x, y)] }$. Definisikan $f$ dengan
\[
f(x) \simeq \umin{y}{R(x, y)}.
\]
Maka $f$ dapat dihitung secara parsial, dan $S$ merupakan domain~$f$.
\end{proof}


\end{document}
